\chapter{Về Thất Bại}
\markboth{Về Thất Bại}{About Failure}

\section{Ai Cũng Thất Bại}
\begin{truyen}{Ai Cũng Thất Bại}{Everyone Fails Sometimes}
\chuhoa{T}{hất bại không phải chuyện chỉ xảy ra với người yếu.}
\textit{(Failure is not something that happens only to weak people.)}

Ai rồi cũng có lúc làm không được, hiểu không ra, hoặc đi sai một đoạn.
\textit{(Everyone has moments of not succeeding, not understanding, or going the wrong way.)}

Biết điều này giúp học sinh đỡ cô đơn hơn khi mình vấp.
\textit{(Knowing this helps students feel less alone when they stumble.)}

Thất bại là phần rất người của việc học và lớn lên.
\textit{(Failure is a very human part of learning and growing.)}
\end{truyen}

\begin{baihoc}
Thất bại không làm em trở thành ngoại lệ; nó chỉ cho thấy em cũng đang đi qua một con đường rất người.
\textit{(Failure does not make you an exception; it shows you are walking a very human road.)}
\end{baihoc}

\begin{ghinhoanh}
\textbf{Short English to remember:}

Failure is common in human growth.

\end{ghinhoanh}

\begin{tuhocnhanh}
1. Vì sao ai cũng có lúc thất bại? \textit{Why does everyone fail at times?}

2. Biết điều này giúp gì cho học sinh? \textit{How does this help students?}

3. Em có đang nghĩ mình cô đơn trong thất bại không? \textit{Do you feel alone in failure?}
\end{tuhocnhanh}
\ngancach

\section{Sai Giúp Mình Học}
\begin{truyen}{Sai Giúp Mình Học}{Mistakes Help Us Learn}
\chuhoa{M}{ột lỗi sai nếu nhìn kỹ thường dạy nhiều hơn một lần làm đúng quá dễ.}
\textit{(A mistake, if examined carefully, often teaches more than an answer that was too easy.)}

Nó chỉ cho em chỗ mình chưa chắc hoặc điều mình đang hiểu lệch.
\textit{(It shows where you are not solid or where your understanding is off.)}

Khi học từ sai, em đang biến điều làm mình ngại thành điều giúp mình khá hơn.
\textit{(When you learn from a mistake, you turn embarrassment into growth.)}

Đó là một khả năng rất quan trọng.
\textit{(That is a very important ability.)}
\end{truyen}

\begin{baihoc}
Sai chỉ thật sự lãng phí khi em không chịu học điều nó đang chỉ ra.
\textit{(A mistake is truly wasted only when you refuse to learn what it is pointing out.)}
\end{baihoc}

\begin{ghinhoanh}
\textbf{Short English to remember:}

Good reflection turns mistakes into lessons.

\end{ghinhoanh}

\begin{tuhocnhanh}
1. Vì sao sai có thể giúp mình học? \textit{Why can mistakes help us learn?}

2. Một lỗi sai thường chỉ điều gì? \textit{What does a mistake usually point to?}

3. Em có lỗi nào đang nên học từ nó? \textit{Is there a mistake you should learn from now?}
\end{tuhocnhanh}
\ngancach

\section{Thất Bại Không Phải Hết}
\begin{truyen}{Thất Bại Không Phải Hết}{Failure Is Not the End}
\chuhoa{M}{ột lần ngã không phải là toàn bộ bản đồ của cuộc đời em.}
\textit{(One fall is not the map of your whole life.)}

Nhiều học sinh sau một lần kết quả không tốt đã vội nghĩ rằng mình hết cơ hội.
\textit{(Many students quickly think they have no chance left after one poor result.)}

Nhưng thất bại thường chỉ là một đoạn khó, không phải điểm cuối.
\textit{(But failure is often only a hard stretch, not the final point.)}

Hiểu điều này giúp em còn giữ được hy vọng để đi tiếp.
\textit{(Understanding this helps you keep enough hope to continue.)}
\end{truyen}

\begin{baihoc}
Đừng dùng một lần chưa được để kết luận quá lớn về tương lai mình.
\textit{(Do not use one bad result to make an enormous conclusion about your future.)}
\end{baihoc}

\begin{ghinhoanh}
\textbf{Short English to remember:}

One failure is not a final identity.

\end{ghinhoanh}

\begin{tuhocnhanh}
1. Vì sao thất bại không phải là hết? \textit{Why is failure not the end?}

2. Nó thường chỉ là điều gì? \textit{What is it often only?}

3. Em có đang kết luận quá sớm về mình không? \textit{Are you concluding too early about yourself?}
\end{tuhocnhanh}
\ngancach

\section{Thử Lại Là Điều Bình Thường}
\begin{truyen}{Thử Lại Là Điều Bình Thường}{Trying Again Is Normal}
\chuhoa{L}{àm lại, bắt đầu lại, và thử lại không phải là điều đáng xấu hổ.}
\textit{(Redoing, restarting, and trying again are not shameful.)}

Có những thứ chỉ hiểu ra sau lần thứ hai, thứ ba, hoặc nhiều hơn.
\textit{(Some things become clear only on the second, third, or later attempt.)}

Người chịu thử lại cho bản thân thêm cơ hội.
\textit{(Those who try again give themselves another chance.)}

Vì vậy, học sinh nên hiểu sớm rằng thử lại là chuyện rất bình thường.
\textit{(So students should understand early that trying again is completely normal.)}
\end{truyen}

\begin{baihoc}
Thử lại không làm em nhỏ đi; nó cho thấy em chưa bỏ mình.
\textit{(Trying again does not make you smaller; it shows you have not abandoned yourself.)}
\end{baihoc}

\begin{ghinhoanh}
\textbf{Short English to remember:}

Trying again keeps hope active.

\end{ghinhoanh}

\begin{tuhocnhanh}
1. Vì sao thử lại là bình thường? \textit{Why is trying again normal?}

2. Người chịu thử lại giữ được điều gì? \textit{What does a person keep by trying again?}

3. Em có việc gì nên thử lại không? \textit{Is there something you should try again?}
\end{tuhocnhanh}
\ngancach

\section{Người Giỏi Cũng Sai}
\begin{truyen}{Người Giỏi Cũng Sai}{Good Students Also Make Mistakes}
\chuhoa{N}{hiều học sinh nhìn người giỏi rồi tưởng họ ít khi sai.}
\textit{(Many students look at strong students and think they rarely make mistakes.)}

Thật ra, người giỏi cũng sai, chỉ là họ học từ sai nhanh hơn và bền hơn.
\textit{(In truth, good students also make mistakes; they just learn from them more steadily.)}

Hiểu điều này giúp em bớt thần tượng hóa người khác và bớt khắt khe với mình.
\textit{(Understanding this helps you stop idealizing others and being too harsh on yourself.)}

Sai không làm mất khả năng của một người.
\textit{(Making mistakes does not erase a person's ability.)}
\end{truyen}

\begin{baihoc}
Người giỏi không phải người không sai, mà là người biết đứng dậy từ chỗ sai.
\textit{(A good learner is not someone who never errs, but someone who can rise from error.)}
\end{baihoc}

\begin{ghinhoanh}
\textbf{Short English to remember:}

Even strong learners still make mistakes.

\end{ghinhoanh}

\begin{tuhocnhanh}
1. Vì sao người giỏi cũng sai? \textit{Why do good students also make mistakes?}

2. Họ khác ở điểm nào? \textit{How are they different?}

3. Em có đang thần tượng hóa người khác quá mức không? \textit{Are you idealizing others too much?}
\end{tuhocnhanh}
\ngancach

\section{Kiên Nhẫn Sau Thất Bại}
\begin{truyen}{Kiên Nhẫn Sau Thất Bại}{Be Patient After Failure}
\chuhoa{S}{au một lần thất bại, con người thường muốn thấy mình khá lại thật nhanh.}
\textit{(After failure, people often want to recover very quickly.)}

Nhưng nhiều vết hụt cần thời gian để hiểu và sửa.
\textit{(But many weak spots need time to understand and correct.)}

Kiên nhẫn sau thất bại giúp em không tự ép mình đến mức kiệt sức hoặc chán nản hơn.
\textit{(Patience after failure keeps you from exhausting yourself or becoming more discouraged.)}

Đi lên sau vấp ngã thường cần thời gian hơn người ta mong.
\textit{(Recovery after stumbling often takes longer than expected.)}
\end{truyen}

\begin{baihoc}
Sau thất bại, vội vàng quá mức không luôn giúp đi nhanh hơn; nhiều khi nó chỉ làm mình mệt thêm.
\textit{(After failure, too much haste does not always help; sometimes it only adds more fatigue.)}
\end{baihoc}

\begin{ghinhoanh}
\textbf{Short English to remember:}

Patience helps recovery after failure.

\end{ghinhoanh}

\begin{tuhocnhanh}
1. Vì sao cần kiên nhẫn sau thất bại? \textit{Why is patience needed after failure?}

2. Vội vàng quá mức có thể gây gì? \textit{What can too much haste cause?}

3. Em có đang đòi mình phục hồi quá nhanh không? \textit{Are you demanding recovery too quickly from yourself?}
\end{tuhocnhanh}
\ngancach

\section{Học Từ Sai Lầm}
\begin{truyen}{Học Từ Sai Lầm}{Learn from Failure}
\chuhoa{S}{ai lầm nào cũng mang theo một bài học nếu em chịu dừng lại để nhìn.}
\textit{(Every failure carries a lesson if you are willing to stop and look.)}

Có khi bài học là về kiến thức, có khi là về thái độ, có khi là về cách chuẩn bị.
\textit{(Sometimes the lesson is about knowledge, sometimes about attitude, and sometimes about preparation.)}

Học từ sai lầm làm thất bại không còn vô nghĩa.
\textit{(Learning from failure keeps it from becoming meaningless.)}

Đó là cách biến một điều buồn thành một điều có ích.
\textit{(It turns something painful into something useful.)}
\end{truyen}

\begin{baihoc}
Sai lầm có thể là thầy dạy khó chịu nhưng rất thật.
\textit{(Failure can be an unpleasant but truthful teacher.)}
\end{baihoc}

\begin{ghinhoanh}
\textbf{Short English to remember:}

Failure can teach if you pay attention.

\end{ghinhoanh}

\begin{tuhocnhanh}
1. Vì sao nên học từ sai lầm? \textit{Why should you learn from failure?}

2. Sai lầm có thể dạy điều gì? \textit{What can failure teach?}

3. Gần đây em rút ra bài học gì? \textit{What lesson have you drawn recently?}
\end{tuhocnhanh}
\ngancach

\section{Đừng Xấu Hổ Vì Sai}
\begin{truyen}{Đừng Xấu Hổ Vì Sai}{Do Not Be Ashamed of Mistakes}
\chuhoa{X}{ấu hổ quá mức vì sai làm nhiều học sinh chỉ muốn trốn khỏi chỗ mình đang yếu.}
\textit{(Too much shame about mistakes makes many students want only to escape their weak points.)}

Nhưng nếu trốn, em khó học được điều mình đang cần học nhất.
\textit{(But if you run away, you miss what you most need to learn.)}

Thừa nhận sai với sự bình tĩnh giúp em lớn hơn nhiều so với việc giữ sĩ diện.
\textit{(Admitting mistakes calmly helps you grow more than protecting pride.)}

Đừng xấu hổ quá lâu vì điều vốn là một phần của học tập.
\textit{(Do not stay ashamed too long about something that is part of learning.)}
\end{truyen}

\begin{baihoc}
Sai không đáng xấu hổ bằng việc từ chối nhìn và sửa cái sai.
\textit{(The mistake itself is less shameful than refusing to face and correct it.)}
\end{baihoc}

\begin{ghinhoanh}
\textbf{Short English to remember:}

Shame should not block correction.

\end{ghinhoanh}

\begin{tuhocnhanh}
1. Vì sao không nên xấu hổ quá mức vì sai? \textit{Why should you not be overly ashamed of mistakes?}

2. Giữ sĩ diện quá mức làm mất điều gì? \textit{What is lost when pride is protected too much?}

3. Em có chỗ sai nào cần nhìn bình tĩnh hơn không? \textit{Is there a mistake you need to face more calmly?}
\end{tuhocnhanh}
\ngancach

\section{Sai Hôm Nay, Đúng Ngày Mai}
\begin{truyen}{Sai Hôm Nay, Đúng Ngày Mai}{Wrong Today, Right Tomorrow}
\chuhoa{C}{ó những điều hôm nay em làm chưa được, nhưng ngày mai quay lại đã khá hơn.}
\textit{(Some things you fail today may become clearer tomorrow.)}

Đầu óc con người cần thời gian để tiêu hóa, sắp xếp và nối lại nhiều mảnh.
\textit{(The human mind needs time to digest, organize, and reconnect many pieces.)}

Vì thế, một lần sai hôm nay chưa quyết định tương lai của em với bài ấy.
\textit{(So one mistake today does not decide your future with that lesson.)}

Chỉ cần em còn quay lại, cánh cửa vẫn chưa khép.
\textit{(As long as you return, the door is not yet closed.)}
\end{truyen}

\begin{baihoc}
Hôm nay chưa đúng không có nghĩa là em mãi mãi không làm được.
\textit{(Not getting it right today does not mean you never will.)}
\end{baihoc}

\begin{ghinhoanh}
\textbf{Short English to remember:}

Today's mistake may become tomorrow's understanding.

\end{ghinhoanh}

\begin{tuhocnhanh}
1. Vì sao sai hôm nay chưa phải là hết? \textit{Why is being wrong today not the end?}

2. Đầu óc cần gì để hiểu rõ hơn? \textit{What does the mind need to understand better?}

3. Em có bài nào nên quay lại vào ngày mai? \textit{What lesson should you revisit tomorrow?}
\end{tuhocnhanh}
\ngancach

\section{Không Bỏ Cuộc}
\begin{truyen}{Không Bỏ Cuộc}{Do Not Give Up}
\chuhoa{Đ}{iều cuối cùng học sinh cần hiểu về thất bại là đừng bỏ cuộc quá sớm.}
\textit{(The final thing students should understand about failure is this: do not give up too early.)}

Không phải vì mọi chuyện sẽ dễ, mà vì nếu bỏ quá nhanh thì gần như chắc chắn không còn cơ hội để đổi khác.
\textit{(Not because everything will be easy, but because quitting too early almost guarantees no chance of change.)}

Đi tiếp qua một thất bại là một bài học rất quan trọng của trưởng thành.
\textit{(Continuing after failure is an important lesson of maturity.)}

Người không bỏ cuộc thường tìm thấy một con đường khác để đi.
\textit{(Those who do not give up often find another way forward.)}
\end{truyen}

\begin{baihoc}
Thất bại có thể làm em chậm lại, nhưng đừng để nó làm em dừng luôn.
\textit{(Failure may slow you down, but do not let it stop you completely.)}
\end{baihoc}

\begin{ghinhoanh}
\textbf{Short English to remember:}

Keep going after failure.

\end{ghinhoanh}

\begin{tuhocnhanh}
1. Vì sao không bỏ cuộc là điều rất quan trọng? \textit{Why is not giving up so important?}

2. Bỏ cuộc quá sớm làm mất gì? \textit{What does quitting too early take away?}

3. Em cần nhắc mình câu nào khi muốn bỏ? \textit{What should you tell yourself when you want to quit?}
\end{tuhocnhanh}
\ngancach